% GENERATED FILE. Edit canonical lessons, then run npm run generate:books.
% canonical-source-hash: fnv1a64:6ac3fe7427d5137b
% canonical-lessons: TA-C03-eppadi, TA-C03-eppadi-irukkirirgal, TA-C03-naan, TA-C03-nalam, TA-C03-paravayillai, TA-C03-practice

\chapter{How Are You?}
\label{ch:responding}
\hlchaptermodality{3}

\begin{chapteropening}
\textbf{By the end of this chapter:} \emph{I can ask how someone is in Tamil, answer naturally, and return the question.}

Take the how-are-you exchange both ways out loud: ask it respectfully, answer that you are well, thank them, and wave the thanks away.
\end{chapteropening}

\section[eppaḍi]{\ta{எப்படி} (eppaḍi) — \textquotedblleft{}how,\textquotedblright{} another of the e- questions}

\label{lesson:TA-C03-eppadi}



\begin{quote}
You met \emph{\textbf{eṉṉa}} (\textquotedblleft{}what\textquotedblright{}) in Chapter 2. Here is its sibling — and, like all of Tamil's question-words, it begins with \textbf{e-}.
\end{quote}

\begin{cousinweb}
\textbf{\ta{எப்படி}} (\emph{eppaḍi}, \textquotedblleft{}how\textquotedblright{}) is native Dravidian, built on the interrogative base \textbf{\ta{எ}-} (\emph{e-}) + \emph{paḍi} (\textquotedblleft{}manner, way\textquotedblright{}) — literally \textquotedblleft{}in what manner.\textquotedblright{} Notice the pattern: Tamil's questions nearly all start with this \textbf{e-} — \emph{eṉṉa} (what), \emph{eppaḍi} (how), \emph{eppōdu} (when), \emph{eṅgē} (where), \emph{ēṉ} (why), and \emph{yār} (who). Where Hindi and English gather their questions under \emph{k-} / \emph{wh-}, Tamil gathers its own under \textbf{e-} — a separate, older family, evolved on its own.
\end{cousinweb}

\subsection*{Your turn}
Say these aloud:

\begin{itemize}
\raggedright
  \item \textquotedblleft{}eppaḍi\textquotedblright{}
  \item the e- question family — \emph{eṉṉa, eppaḍi, eppōdu, eṅgē, ēṉ}
  \item what it's built from (\emph{e-} \textquotedblleft{}what\textquotedblright{} + \emph{paḍi} \textquotedblleft{}manner\textquotedblright{})
\end{itemize}

\begin{tcolorbox}[breakable,colback=teal!4,colframe=teal!35!black,title={Before you move on}]
What base do Tamil's question-words share, and how is it different from Hindi's? (The interrogative \emph{e-}; it is native Dravidian, unrelated to Hindi's \emph{k-} / English's \emph{wh-}.)
\end{tcolorbox}
\section[nīṅgaḷ eppaḍi irukkiṟīrgaḷ?]{\ta{நீங்கள்} \ta{எப்படி} \ta{இருக்கிறீர்கள்}? (nīṅgaḷ eppaḍi irukkiṟīrgaḷ?)}

\label{lesson:TA-C03-eppadi-irukkirirgal}



\begin{quote}
The everyday \textquotedblleft{}how are you?\textquotedblright{} — and the first time Tamil actually \emph{needs} a verb \textquotedblleft{}to be.\textquotedblright{}
\end{quote}

\begin{cousinweb}
\textbf{\ta{நீங்கள்} \ta{எப்படி} \ta{இருக்கிறீர்கள்}?} = \textbf{\ta{நீங்கள்}} (\emph{nīṅgaḷ}, \textquotedblleft{}you,\textquotedblright{} respectful) + \textbf{\ta{எப்படி}} (\emph{eppaḍi}, \textquotedblleft{}how\textquotedblright{}) + \textbf{\ta{இருக்கிறீர்கள்}} (\emph{irukkiṟīrgaḷ}, \textquotedblleft{}are/exist\textquotedblright{}) — \textquotedblleft{}you how are?\textquotedblright{} The verb comes \textbf{last}, as in every Tamil sentence.

The verb is \textbf{\ta{இரு}} (\emph{iru}), \textquotedblleft{}to be, to exist, to stay\textquotedblright{} — native Dravidian. Inside \emph{irukkiṟīrgaḷ} that stem carries two endings: \emph{iru} + \emph{-kkiṟ-} (present) + \emph{-īrgaḷ} (\textquotedblleft{}you,\textquotedblright{} respectful/plural) — the three run together exactly as written. A long word, but built of clean pieces, and you can hear all three.
\end{cousinweb}

\begin{grammarlens}[title={Grammar Lens: now the copula returns}]
In Chapter 2 you learned Tamil drops \textquotedblleft{}is\textquotedblright{} entirely: \emph{eṉ peyar Arun} (\textquotedblleft{}my name Arun\textquotedblright{}). That \textbf{zero copula} works for \emph{equations} (X is Y). But \textquotedblleft{}how are you?\textquotedblright{} is about a \textbf{state} — how you \emph{exist} right now — and for that Tamil does use a verb: \textbf{\ta{இரு}} (\emph{iru}). So: no verb for \textquotedblleft{}my name is Arun,\textquotedblright{} but a real verb for \textquotedblleft{}how are you.\textquotedblright{} The verb ending \textbf{-\ta{ீர்கள்}} (\emph{-īrgaḷ}) already means \textquotedblleft{}you (respectful),\textquotedblright{} so \emph{nīṅgaḷ} can even be dropped.
\end{grammarlens}

\subsection*{Your turn}
Say these aloud:

\begin{itemize}
\raggedright
  \item the respectful question — \emph{nīṅgaḷ eppaḍi irukkiṟīrgaḷ?}
  \item the verb and its meaning (\emph{iru} — \textquotedblleft{}to be / exist / stay\textquotedblright{})
  \item why here Tamil needs a verb but \textquotedblleft{}my name is Arun\textquotedblright{} doesn't (state vs. equation)
\end{itemize}

\begin{tcolorbox}[breakable,colback=teal!4,colframe=teal!35!black,title={Before you move on}]
What verb does \textquotedblleft{}how are you?\textquotedblright{} use, and why does this sentence need one when \textquotedblleft{}my name is Arun\textquotedblright{} didn't? (\emph{Iru}, \textquotedblleft{}to be/exist\textquotedblright{}; because it describes a state, not an equation — the zero copula only covers X is Y.)
\end{tcolorbox}
\section[nāṉ]{\ta{நான்} (nāṉ) — \textquotedblleft{}I,\textquotedblright{} the Dravidian first person}

\label{lesson:TA-C03-naan}



\begin{quote}
To answer \textquotedblleft{}how are you?\textquotedblright{} you first need \textquotedblleft{}I.\textquotedblright{} In Chapter 2 you learned \emph{eṉ} (\textquotedblleft{}my\textquotedblright{}); here is its owner.
\end{quote}

\begin{cousinweb}
\textbf{\ta{நான்}} (\emph{nāṉ}, \textquotedblleft{}I\textquotedblright{}) $\leftarrow$ Proto-Dravidian \textbf{*yāṉ / *nāṉ} — a native first-person word, and (unlike Hindi's \emph{maiṁ}, cousin of English \emph{me}) \textbf{not} related to the European \textquotedblleft{}I/me\textquotedblright{} family at all. Its possessive is the \emph{\textbf{eṉ}} (\textquotedblleft{}my\textquotedblright{}) you already met: \emph{nāṉ} \textquotedblleft{}I,\textquotedblright{} \emph{eṉ} \textquotedblleft{}my\textquotedblright{} — the same root, one standing, one leaning on a noun. Tamil's \textquotedblleft{}I\textquotedblright{} is one of the clearest markers that Dravidian grew up entirely apart from the Indo-European world around it.
\end{cousinweb}

\begin{grammarlens}[title={Grammar Lens: \textquotedblleft{}I\textquotedblright{} you can drop}]
Tamil verbs carry the person in their ending — \emph{irukk-iṟēṉ} already means \textquotedblleft{}\textbf{I} am\textquotedblright{} — so \emph{nāṉ} is often left out once it's clear. You will still learn it first, because the answer to \emph{eppaḍi irukkiṟīrgaḷ?} begins with it.
\end{grammarlens}

\subsection*{Your turn}
Say these aloud:

\begin{itemize}
\raggedright
  \item \textquotedblleft{}nāṉ\textquotedblright{}
  \item the pair — \emph{nāṉ} (I), \emph{eṉ} (my)
  \item is it related to English \emph{me}? (No — Dravidian, a separate family)
\end{itemize}

\begin{tcolorbox}[breakable,colback=teal!4,colframe=teal!35!black,title={Before you move on}]
What is \emph{nāṉ}'s possessive, and is it cousin to English \emph{me}? (\emph{Eṉ}, \textquotedblleft{}my\textquotedblright{}; no — it is native Dravidian, unrelated to the Indo-European pronouns.)
\end{tcolorbox}
\section[nalam]{\ta{நலம்} (nalam) — \textquotedblleft{}well,\textquotedblright{} and how to answer}

\label{lesson:TA-C03-nalam}



\begin{quote}
The word that answers \textquotedblleft{}how are you?\textquotedblright{} — and a whole family of Tamil \textquotedblleft{}good\textquotedblright{} grows from it.
\end{quote}

\begin{cousinweb}
\textbf{\ta{நலம்}} (\emph{nalam}, \textquotedblleft{}wellness, good, welfare\textquotedblright{}) is native Dravidian, from the root \emph{nal-} (\textquotedblleft{}good\textquotedblright{}) — the same root behind \emph{nalla} (\textquotedblleft{}good,\textquotedblright{} an everyday adjective) and \emph{naṉṟi} (\textquotedblleft{}thanks\textquotedblright{} — literally \textquotedblleft{}goodness,\textquotedblright{} which you met in Chapter 1). So \textquotedblleft{}thank you\textquotedblright{} and \textquotedblleft{}I'm well\textquotedblright{} in Tamil share a root: both are \emph{nal-}, \textquotedblleft{}good.\textquotedblright{} A purely Dravidian word, with no Sanskrit or European cousin — Tamil answering in its own voice.
\end{cousinweb}

\subsection*{You'll want to know: \ta{நலம்}}
Assemble it: \emph{\textbf{nāṉ nalamāga irukkiṟēṉ}} — \textquotedblleft{}\textbf{I am well}.\textquotedblright{} The piece \emph{nalam} + \emph{-āga} (\textquotedblleft{}in a … way\textquotedblright{}) makes the adverb \emph{nalamāga} (\textquotedblleft{}well-ly\textquotedblright{}); \emph{irukkiṟēṉ} is \textquotedblleft{}I am/exist\textquotedblright{} (the \emph{iru} verb, first person). The whole exchange:

\begin{quote}
— \emph{nīṅgaḷ eppaḍi irukkiṟīrgaḷ?} (\textquotedblleft{}How are you?\textquotedblright{}) — \emph{nāṉ nalamāga irukkiṟēṉ, naṉṟi.} (\textquotedblleft{}I'm well, thank you.\textquotedblright{})
\end{quote}

Casually, Tamils shorten it to \emph{nallā irukkēṉ} (\textquotedblleft{}I'm good\textquotedblright{}).

\subsection*{Your turn}
Say these aloud:

\begin{itemize}
\raggedright
  \item \textquotedblleft{}nalam\textquotedblright{}
  \item the full reply — \emph{nāṉ nalamāga irukkiṟēṉ}
  \item the root it shares with \textquotedblleft{}thanks\textquotedblright{} (\emph{nal-} \textquotedblleft{}good,\textquotedblright{} as in \emph{naṉṟi})
\end{itemize}

\begin{tcolorbox}[breakable,colback=teal!4,colframe=teal!35!black,title={Before you move on}]
Give the full reply to \textquotedblleft{}how are you?\textquotedblright{}, and name the Chapter-1 word that shares \emph{nalam}'s root. (\emph{Nāṉ nalamāga irukkiṟēṉ}; \emph{naṉṟi}, \textquotedblleft{}thanks\textquotedblright{} — both from \emph{nal-}, \textquotedblleft{}good.\textquotedblright{})
\end{tcolorbox}
\section[paravāyillai]{\ta{பரவாயில்லை} (paravāyillai) — \textquotedblleft{}no problem\textquotedblright{}}

\label{lesson:TA-C03-paravayillai}



\begin{quote}
When someone thanks you — or apologises — this is the easy, gracious reply, and it hides one of the most Tamil words there is.
\end{quote}

\begin{cousinweb}
\textbf{\ta{பரவாயில்லை}} literally means \textquotedblleft{}\textbf{there is no harm}\textquotedblright{} — \emph{paravāy} (\textquotedblleft{}harm, trouble\textquotedblright{}) + \emph{illai} (\textquotedblleft{}is-not, there-is-not\textquotedblright{}). It serves as \textquotedblleft{}it's okay / no worries / never mind / you're welcome.\textquotedblright{} The key atom, \emph{illai}, you already met in Chapter 1 — the Tamil word for negation and non-existence, the mirror of \emph{iru} (\textquotedblleft{}to be\textquotedblright{}). Where \emph{iru} says \textquotedblleft{}it is,\textquotedblright{} \emph{illai} says \textquotedblleft{}it is not.\textquotedblright{} This one word is shared, in almost the same shape, across Tamil, Kannada, and Malayalam (\emph{illa}) — one of the deep bonds of the Dravidian south. (Telugu alone breaks away, using \emph{lēdu}.)
\end{cousinweb}

\begin{grammarlens}[title={Grammar Lens: \emph{illai}, the all-purpose \textquotedblleft{}no\textquotedblright{}}]
\emph{Illai} negates existence and equations alike: \emph{nāṉ illai} (\textquotedblleft{}it's not me\textquotedblright{}), \emph{paṇam illai} (\textquotedblleft{}there's no money\textquotedblright{}). It is also, on its own, one way to say plain \textquotedblleft{}no.\textquotedblright{} Learn \emph{iru} / \emph{illai} together — \textquotedblleft{}is\textquotedblright{} and \textquotedblleft{}is-not\textquotedblright{} — and you can already affirm or deny almost anything.
\end{grammarlens}

\subsection*{Your turn}
Say these aloud:

\begin{itemize}
\raggedright
  \item \textquotedblleft{}paravāyillai\textquotedblright{}
  \item what it literally means (\textquotedblleft{}there is no harm\textquotedblright{})
  \item the \emph{iru} / \emph{illai} pair — \textquotedblleft{}is\textquotedblright{} and \textquotedblleft{}is-not\textquotedblright{}
\end{itemize}

\begin{tcolorbox}[breakable,colback=teal!4,colframe=teal!35!black,title={Before you move on}]
What does \emph{paravāyillai} literally say, and which Dravidian languages share its \emph{illai}? (\textquotedblleft{}There is no harm\textquotedblright{}; Tamil, Kannada, Malayalam — Telugu uses \emph{lēdu} instead.)
\end{tcolorbox}
\section[Practice]{Chapter 3 — The how-are-you exchange}

\label{lesson:TA-C03-practice}



\begin{quote}
No new words. The whole responding cycle, from atoms you now own.
\end{quote}

\subsection*{The exchange}
Say these aloud:

\begin{itemize}
\raggedright
  \item \textquotedblleft{}How are you?\textquotedblright{} (respectful) — \emph{nīṅgaḷ eppaḍi irukkiṟīrgaḷ?}
  \item \textquotedblleft{}I'm well, thank you.\textquotedblright{} — \emph{nāṉ nalamāga irukkiṟēṉ, naṉṟi.}
  \item \textquotedblleft{}No problem / you're welcome.\textquotedblright{} — \emph{paravāyillai.}
  \item the casual version — \emph{nallā irukkēṉ.}
\end{itemize}

\subsection*{Your turn — the atoms}
Say these aloud:

\begin{itemize}
\raggedright
  \item the e- question-word for \textquotedblleft{}how\textquotedblright{} (\emph{eppaḍi})
  \item \textquotedblleft{}I\textquotedblright{} and the verb \textquotedblleft{}to be\textquotedblright{} (\emph{nāṉ}, \emph{iru})
  \item the \textquotedblleft{}is / is-not\textquotedblright{} pair (\emph{iru} / \emph{illai})
\end{itemize}

\begin{grammarlens}[title={What you've built — roots you now carry}]
\begin{itemize}
\raggedright
  \item \emph{eppaḍi} $\leftarrow$ the interrogative \emph{e-} (native Dravidian, not Hindi's \emph{k-})
  \item \emph{nāṉ} $\leftarrow$ Proto-Dravidian \emph{*nāṉ} (unrelated to English \emph{me}); possessive \emph{eṉ}
  \item \emph{nalam} $\leftarrow$ \emph{nal-} \textquotedblleft{}good\textquotedblright{} — the same root as \emph{naṉṟi} (\textquotedblleft{}thanks\textquotedblright{})
  \item \emph{iru} / \emph{illai} — \textquotedblleft{}to be\textquotedblright{} and \textquotedblleft{}to be-not\textquotedblright{}; \emph{illai} shared with Kannada \& Malayalam
\end{itemize}
\end{grammarlens}

\begin{grammarlens}[title={Grammar lens — what you now hold}]
\begin{itemize}
\raggedright
  \item \emph{Say it:} when Tamil needs the verb \emph{iru} and when it drops the copula (state vs. equation)
\end{itemize}
\end{grammarlens}

\begin{tcolorbox}[breakable,colback=teal!4,colframe=teal!35!black,title={Before you move on}]
A stranger asks \emph{nīṅgaḷ eppaḍi irukkiṟīrgaḷ?} — give a full, polite reply, and answer their thanks. (\emph{Nāṉ nalamāga irukkiṟēṉ, naṉṟi} — and to thanks, \emph{paravāyillai}.)
\end{tcolorbox}
